\section{Optimization problem}

\subsection{Linear decomposition of the solution}

By linearity of equation~\eqref{eq:VF}, we decompose $f = \sum_{i=1}^6 \alpha_i \mathbf{1}_{C_i}$ and write:
\[
T(\alpha, x) = T_0(x) + \sum_{i=1}^6 \alpha_i\,T_i(x),
\]
where $T_0$ is the solution of the problem \emph{without source} ($f=0$) with the same boundary conditions (Dirichlet and Robin), and for $i=1,\ldots,6$, $T_i$ is the solution of the homogeneous Dirichlet temperature problem ($T_i|_{\Gamma_{\text{haut}}}=0$) with the right-hand side $f = \mathbf{1}_{C_i}$:
\[
-\mathrm{div}(k\nabla T_i) = \mathbf{1}_{C_i} \quad \text{in } \Omega, \quad
T_i|_{\Gamma_{\text{haut}}} = 0, \quad
\partial_n T_i + h T_i = 0 \text{ on } \Gamma_{\text{bas}}, \quad
\partial_n T_i = 0 \text{ on } \Gamma_{\text{lat}}.
\]
These 7 problems are independent of the coefficients $(\alpha_i)$ and only need to be solved once.

\subsection{Reduction to a linear system}

We seek $\alpha^* = (\alpha_i^*)_{i=1}^6$ minimizing the functional:
\[
J(\alpha) = \frac{1}{2}\int_S |T(\alpha,x) - T_s|^2\dd x.
\]
By substituting the decomposition, $J$ becomes a convex quadratic function in $\alpha$:
\[
J(\alpha) = \frac{1}{2}\int_S \left|T_0 + \sum_{i=1}^6 \alpha_i T_i - T_s\right|^2\dd x.
\]
The optimality conditions $\partial J/\partial \alpha_j = 0$ (for $j=1,\ldots,6$) give:
\[
\int_S \left(T_0 + \sum_{i=1}^6 \alpha_i T_i - T_s\right)T_j\dd x = 0, \quad j = 1,\ldots,6,
\]
which yields the $6\times 6$ linear system:
\begin{equation}\label{eq:syst_opt}
\boxed{M\,\alpha^* = b, \qquad M_{ij} = \int_S T_i\,T_j\dd x, \qquad b_j = \int_S (T_s - T_0)\,T_j\dd x.}
\end{equation}
The matrix $M$ is symmetric positive semi-definite. It is positive definite if the functions $(T_i|_S)_{i=1}^6$ are linearly independent, which is the generic case.

\subsection{Optimal coefficients for $T_s = 400\,\mathrm{K}$}

System~\eqref{eq:syst_opt} is solved by Gaussian elimination with partial pivoting (implemented manually in the FreeFem++ code). The solution $T^{\text{opt}} = T_0 + \sum_i \alpha_i^* T_i$ is then compared to a direct simulation using these coefficients, which validates the procedure. The quality is measured by the standard deviation $\sigma_S = \left(\frac{1}{|S|}\int_S (T^{\text{opt}} - T_s)^2\dd x\right)^{1/2}$ and by the value of $J(\alpha^*)$.

\subsection{Improvement of the distribution by adding sources}

By increasing the number of resistors (going to 8 sources in the code, placed on a grid $\{-1, -1/3, 1/3, 1\} \times \{-0.75, 0.75\}$), the same linear formalism applies with an $8\times8$ system. By altering the layout and shape of the sources, we can significantly reduce the temperature standard deviation in $S$, at the cost of a slight increase in power requirements (higher intensities). Sources closer to the center of $S$ are more efficient, but since the resistors are placed outside of $S$, a geometric compromise is necessary.
